مشکل اصلی در این سناریو، {تجزیه زودرس سیستم} یا همان \lr{Premature Decomposition} (یا اصطلاحاً \lr{Distributed Monolith}) است. تیم توسعه پیش از درک کامل بیزینس و مرزهای دامنه مسئله (\lr{Bounded Contexts})، سیستم را به میکروسرویس‌های متعدد تقسیم کرده است. این کار بدون وجود نیاز واقعی به مقیاس‌پذیری، تنها پیچیدگی‌های عملیاتی را به شدت افزایش داده است.

شروع مستقیم با معماری میکروسرویس در ابتدای پروژه به دلایل زیر اشتباه است:
\begin{itemize}
    \item {عدم شناخت دامنه:} وقتی مرزهای بیزینس شفاف نیست، میکروسرویس‌ها به اشتباه تفکیک می‌شوند. این امر منجر به ایجاد \lr{Domain Coupling} شدید شده و تغییر در یک ویژگی، نیازمند تغییر همزمان در چندین سرویس و \lr{API} می‌شود.
    \item {ایجاد منولیت توزیع‌شده:} سیستم تمام معایب منولیت (وابستگی شدید) و تمام معایب میکروسرویس (پیچیدگی شبکه، دشواری تست و باگ‌های مرزی) را یکجا دارد.
    \item {کاهش سرعت توسعه:} نوشتن تست‌های \lr{End-to-End} در یک سیستم توزیع‌شده با دیتابیس‌های مجزا بسیار سخت‌تر از یک سیستم یکپارچه است که نتیجه آن افت شدید سرعت تحویل \lr{Feature}هاست.
\end{itemize}

رویکرد استاندارد و پیشنهادی، استراتژی {\lr{Monolith First}} است:
\begin{enumerate}
    \item {شروع با یک منولیت ماژولار:} سیستم باید در ابتدا به صورت یکپارچه اما با کدهای ماژولار و ساختاریافته توسعه یابد. در این حالت، مرزها در سطح کد (پکیج‌ها یا کامپوننت‌ها) تعریف می‌شوند و ارتباطات به جای شبکه، از طریق فراخوانی توابع داخل حافظه است.
    \item {کشف مرزهای دامنه:} با رشد پروژه و شناخت دقیق‌تر رفتار کاربران و بیزینس، مرزهای واقعی دامنه‌ها (\lr{Subdomains}) مشخص می‌شوند.
    \item {مهاجرت تدریجی در صورت نیاز:} هر زمان که یک ماژول خاص از نظر فنی نیاز به مقیاس‌پذیری مستقل پیدا کرد یا تیم‌های توسعه آن‌قدر بزرگ شدند که روی یک کدبیس تداخل داشتند، آن ماژول به صورت یک میکروسرویس مجزا جدا می‌شود.
\end{enumerate}
